% !TeX root = ../main.tex
% Original qualitative counterfactual examples, with all other conditions fixed.
% Row 1: same geometry, opposite initial velocity, same upward action.
% Row 2: same body geometry, connected/separated pair, same action on the right body.
% Pale outlines mark current reference positions; dashed outlines are conceptual
% future responses, not measured trajectories or actual model predictions.
\begingroup
% Shared local styles for the four original problem illustrations in Section 2.2.
% Loaded inside each figure's group; all meanings remain legible in grayscale.
\linespread{1}\selectfont
\tikzset{
  rpBase/.style={x=1cm,y=1cm,text=black,line cap=round,line join=round,
    font=\songti\fontsize{9.8pt}{12pt}\selectfont,
    every node/.style={align=center,inner sep=0pt,outer sep=0pt}},
  rpTitle/.style={font=\heiti\fontsize{10.3pt}{13pt}\selectfont},
  rpNote/.style={font=\songti\fontsize{9.3pt}{12pt}\selectfont},
  rpFrame/.style={draw=black,line width=0.65pt},
  rpArrow/.style={draw=black,line width=0.8pt,
    -{Latex[length=1.7mm,width=1.15mm]}},
  rpMotion/.style={draw=black,line width=0.95pt,
    -{Latex[length=1.6mm,width=1.1mm]}},
  rpAction/.style={draw=black,line width=1.5pt,
    -{Latex[length=2.2mm,width=1.7mm]}},
  rpHistory/.style={draw=black!65,line width=0.65pt,
    dash pattern=on 1.4pt off 1.6pt},
  rpFuture/.style={draw=black,line width=0.75pt,
    dash pattern=on 3pt off 2pt},
  rpBoundary/.style={draw=black,line width=1.25pt}
}
\newcommand{\rpPanel}[5]{%
  \draw[rpFrame] (#1,#2) rectangle (#3,#4);
  \path[fill=black!5] (#1+0.01,#4-0.55) rectangle (#3-0.01,#4-0.01);
  \draw[rpFrame] (#1,#4-0.55) -- (#3,#4-0.55);
  \node[rpTitle] at ({(#1+#3)/2},{#4-0.275}) {#5};
}
\newcommand{\rpFooter}[1]{%
  \draw[rpFrame] (0.10,0.10) rectangle (15.90,0.77);
  \node at (8,0.435) {#1};
}
% Hatching is explicitly drawn and clipped, so no additional TikZ libraries
% or color meanings are required. Argument: a closed local-coordinate path.
\newcommand{\rpHatchFace}[1]{%
  \path[fill=black!5] #1;
  \begin{scope}
    \clip #1;
    \foreach \rhx in {-6,-5.8,...,8}{%
      \draw[black,line width=0.35pt] (\rhx,-1) -- ++(3,6);
    }
  \end{scope}
  \draw[rpBoundary] #1;
}

\begin{tikzpicture}[rpBase]
  \path[use as bounding box] (0,0) rectangle (16,6.75);

  \path[draw=black,fill=black!5,line width=0.65pt] (0.10,6.05) rectangle (15.90,6.60);
  \foreach \xx in {2.30,7.45,10.10}{\draw[rpFrame] (\xx,6.05) -- (\xx,6.60);}
  \node[rpTitle] at (1.20,6.325) {状态差异};
  \node[rpTitle] at (4.875,6.325) {当前状态};
  \node[rpTitle] at (8.775,6.325) {同一动作};
  \node[rpTitle] at (13.00,6.325) {未来响应示意};
  \foreach \yb/\yt in {3.55/5.95,1.00/3.40}{%
    \draw[rpFrame] (0.10,\yb) rectangle (15.90,\yt);
    \foreach \xx in {2.30,7.45,10.10}{\draw[rpFrame] (\xx,\yb) -- (\xx,\yt);}
  }

  \node[rpTitle] at (1.20,5.02) {运动趋势};
  \node[rpNote] at (1.20,4.48) {当前几何相同};
  \foreach \yy in {5.12,4.02}{%
    \draw[draw=black!45,line width=0.4pt] (3.00,\yy-0.30) -- (6.70,\yy-0.30);
    \filldraw[draw=black,fill=black!5,line width=0.75pt] (4.60,\yy) circle[radius=0.22];
    \draw[rpAction] (8.775,\yy-0.27) -- (8.775,\yy+0.40);
    \node at (9.36,\yy+0.10) {$a$};
  }
  \draw[rpMotion] (4.30,5.52) -- (3.35,5.52);
  \node at (5.15,5.52) {$v_0$};
  \draw[rpMotion] (4.90,4.42) -- (5.85,4.42);
  \node at (4.10,4.42) {$v_0$};

  \draw[draw=black!45,line width=0.55pt] (13.00,5.12) circle[radius=0.22];
  \draw[draw=black!45,line width=0.55pt] (13.00,4.02) circle[radius=0.22];
  \draw[rpFuture,-{Latex[length=1.5mm,width=1mm]}]
    (12.80,5.15) .. controls (12.40,5.18) and (12.05,5.35) .. (11.78,5.48);
  \draw[rpFuture,-{Latex[length=1.5mm,width=1mm]}]
    (13.20,4.05) .. controls (13.60,4.08) and (13.95,4.25) .. (14.22,4.38);
  \draw[rpFuture] (11.55,5.53) circle[radius=0.22];
  \draw[rpFuture] (14.45,4.43) circle[radius=0.22];

  \node[rpTitle] at (1.20,2.48) {连接关系};
  \node[rpNote] at (1.20,1.94) {物体几何相近};
  \foreach \yy in {2.65,1.55}{%
    \draw[draw=black,fill=black!5,line width=0.65pt] (4.00,\yy-0.24) rectangle ++(0.62,0.48);
    \draw[draw=black,fill=white,line width=0.65pt] (5.35,\yy-0.24) rectangle ++(0.62,0.48);
    \draw[rpAction] (8.15,\yy) -- (9.20,\yy);
    \node at (9.56,\yy) {$a$};
  }
  \draw[draw=black,line width=1.05pt] (4.62,2.65) -- (5.35,2.65);
  \node[rpNote] at (3.07,2.65) {连接};
  \node[rpNote] at (3.07,1.55) {分离};

  % Reference and future positions share a coordinate frame in both cases.
  \foreach \yy in {2.65,1.55}{%
    \draw[draw=black!40,line width=0.5pt] (10.65,\yy-0.24) rectangle ++(0.50,0.48);
    \draw[draw=black!40,line width=0.5pt] (11.80,\yy-0.24) rectangle ++(0.50,0.48);
    \draw[rpFuture] (13.60,\yy-0.24) rectangle ++(0.50,0.48);
  }
  \draw[rpFuture] (12.45,2.41) rectangle ++(0.50,0.48);
  \draw[rpFuture] (12.95,2.65) -- (13.60,2.65);
  \draw[rpFuture] (10.65,1.31) rectangle ++(0.50,0.48);
  \node[rpNote] at (14.95,2.65) {整体响应};
  \node[rpNote] at (14.95,1.55) {局部响应};

  \rpFooter{研究关注：运动趋势和有效连接如何约束动作响应？}
\end{tikzpicture}
\endgroup
